\documentclass[12pt]{article}
\usepackage{xr,geometry,fancyhdr,hyperref,ifpdf,amsmath,rcs}
\usepackage{color,lastpage,longtable,Ventry,url,paunits,shortcuts,smallsec,float}
\geometry{letterpaper,top=50pt,hmargin={20mm,20mm},headheight=15pt} 

\pagestyle{fancy} 


\fancypagestyle{first}{
\lhead{A405}
\chead{Accurate expressions for $h$ and $\theta_e$ }}
\rhead{page~\thepage/\pageref{LastPage}}
\lfoot{} 
\cfoot{} 
\rfoot{}


\ifpdf
    \usepackage[pdftex]{graphicx} 
    \usepackage{hyperref}
    \pdfcompresslevel=0
    \DeclareGraphicsExtensions{.pdf,.jpg,.mps,.png}
\else
    \usepackage{hyperref}
    \usepackage[dvips]{graphicx}
    \DeclareGraphicsRule{.eps.gz}{eps}{.eps.bb}{`gzip -d #1}
    \DeclareGraphicsExtensions{.eps,.eps.gz}
\fi
\externaldocument[lec6,]{lec6}

\begin{document}

\pagestyle{first}


\section{Review:  static energy, entropy and $\theta_e$}

In this section I repeat the derivation in the 
Week 5 \href{https://drive.google.com/file/d/1BPTvFijtktG4135T1X8DXwI8l_C8xHEt/view?usp=sharing}{static energy and $\theta_e$ notes} to to remove two approximations we made then.

1) We approximated the heat capacity for cloudy air with $c_{pd}$, the heat capacity of dry air $c_{pd}$. This introduces an error of a few percent.
2) As noted on p.7 equation 47, and in the
\href{https://phaustin.github.io/a405/notebooks/worksheets/worksheet7_taylor_solution.html}{Worksheet 7 on Taylor Series} we made the approximation:

\begin{equation}
d  \left ( \frac{l_v(T) r_s(T)}{c_p T} \right ) \approx  \left ( \frac{l_v(T) }{c_p T} \right ) dr_s
\end{equation}
which makes and error as large as 30\%.  These errors are tolerable for qualitative arguments, but if you need to integrate a forecast or climate model over thousands of timesteps you need to do much better than that.

Below I'm going to follow the text ``Atmospheric Convection'' by Kerry Emanuel (1994) and derive a new
version of the equivalent potential temperature that works in
subsaturated air and is fully reversible, without making either approximation.


\subsection{Enthalpy}
\label{sec:enthalpy}


First, review what we know about the enthalpy of a parcel with both water vapor and liquid water,
with total water given by $r_t = r_v + r_l$

\begin{equation}
  \label{eq:h1}
  h = h_d + r_v h_v + r_l h_l = c_{pd} T + r_v c_{pv} T + r_l c_l T
\end{equation}
add and subtract $r_v c_l T$:
\begin{equation}
  \label{eq:h1b}
  h = c_{pd} T + r_v c_{pv} T + r_l c_l T + r_v c_l T - r_v c_l T
\end{equation}
Rearrange:

\begin{equation}
  \label{eq:h2}
  h = c_{pd} T + r_v c_l T + r_l c_l T + r_v c_{pv} T -  r_v c_l T 
\end{equation}
Remember that $l_v= (h_v - h_l) = c_{pv} T -   c_l T$, so that 
\begin{equation}
  \label{eq:h3}
  h = c_{pd} T + (r_v   + r_l )c_l\, T  + r_v(h_v - h_l)
\end{equation}
or in other words
\begin{equation}
  \label{eq:h4}
  h   = (c_{pd}  + r_T c_l) T + r_v l_v
\end{equation}


Note that, even though the latent heat $l_v$ appears in \eqref{eq:h4}, I haven't
made any assumptions about conservation of water, and I have not had to
assume saturation.  

It would not be a waste of time to take the
differential of \eqref{eq:h4} and convince yourself that it gives:


\begin{equation}
  \label{eq:h4diff}
  dh = \underbrace{(c_{pd} + r_v c_{pv} + r_l c_l ) dT}_a + 
  \underbrace{c_l T d r_T}_b + \underbrace{l_v d r_v}_c
\end{equation}
We've always assumed $dr_T = 0$ and discarded term $b$, but it's not a
necessary approximation, just convenient.  For example we could have added it
back in when we increased $r_T$ by 
ocean evaporation when we solved the  tropical heat engine problem.

\subsection{First and second law}
\label{sec:first-second-law}

No matter whether it's a pure substance or a mixture, the atmospheric version of the first
and (reversible) second laws will hold:  
\begin{equation}
  \label{eq:first}
 q\,dt = T d\phi = c_p dT - v dp = dh - v dp
\end{equation}


And if the process is irreversible, then
\begin{equation}
  \label{eq:firstirr}
 q\,dt < T d\phi = c_p dT - v dp = dh - v dp
\end{equation}



Thus for the entropy of dry air, vapor and liquid, I can integrate
\eqref{eq:firstirr} from some agreed-upon base state to get:


\begin{subequations}
    \label{eq:entropy}
  \begin{eqnarray}
\phi_d& =& c_{pd} \log T - R_d \log p_d\\
\phi_v &=&c_{pv} \log T - R_v \log e \\
\phi_l &=&c_l \log T
  \end{eqnarray}
\end{subequations}
and just as for the enthalpy equation \eqref{eq:h1}, the entropy of the mixture is:

\begin{equation}
  \label{eq:entmix}
  \phi = \phi_d + r_v \phi_v + r_l \phi_l
\end{equation}

In equation \eqref{eq:h1b} above, I added and subtracted $r_v c_l T$
so that I could get the latent heat out where I could see it.  I want
to do the same thing here, but it's more complicated for the entropy
because I also need to work with the vapor pressure $e$.  From
\href{https://drive.google.com/file/d/1rJO32NGhTW3TA6noQmSdNl2GNXyrZUgA/view}%
{the Maxwell relation notes} we know that in equilibrium $g_l = g_v$, where $g_l= h_l - T\phi_l$
and $g_v= h_v - T \phi_{vs}$.  I'll use the the subscript $_s$ to indicate that we're
talking about the entropy of saturated vapor in equilibrium (given by
the Clausius-Clapeyron equation) as opposed to the actual entropy
$\phi_v$.  So since $g_l=g_v$ in equilibrium I have:

\begin{equation}
  \label{eq:cc}
  h_l  - T\phi_l = h_v  - T \phi_{vs}
\end{equation}
or in other words:

\begin{equation}
  \label{eq:cc2}
 l_v = h_v - h_l = T ( \phi_{vs} - \phi_l)
\end{equation}
which is just a form of the Clausius-Clapeyron equation.

Given this, I use  \eqref{eq:cc2} to add and subtract
$r_v l_v/T$ from \eqref{eq:entmix}
so that the effect of phase change  on the entropy is clearly visible:

\begin{equation}
  \label{eq:entmix2}
  \phi = \phi_d + r_t  \phi_l + \frac{l_v r_v}{T} + r_v( \phi_v - \phi_{vs}) 
\end{equation}
Note that the last term is zero if we're at saturation.


So now put in the definitions from \eqref{eq:entropy} into
\eqref{eq:entmix2} and rearrange.  Convince yourself that the exact
form for the entropy of air/vapor/liquid is then:


\begin{equation}
  \label{eq:sexact}
  \phi = (c_{pd} + r_T c_l) \log T - R_d \log p_d + \frac{l_v r_v}{T} - r_v R_v \log (e/e_s)
\end{equation}


Now add $R_d \log p_0$ to both sides of \eqref{eq:sexact} and write the
heat capacity as $\widehat{c_p} = c_{pd} + r_T c_l$.

\begin{equation}
  \label{eq:sexact2}
  \phi + R_d \log{p_0}= 
\widehat{c_p} \log T - R_d \log \left ( \frac{p_0}{p_d} \right )  
+ \frac{l_v r_v}{T} - r_v R_v \log (e/e_s)
\end{equation}


Finally, divide by $\widehat{c_p}$  and take $\exp$ of both sides and call the result $\theta_e$:

\begin{equation}
  \label{eq:truetheta}
 \theta_e = \exp \left ( \left ( \phi + R_d \log{p_0} \right )/ \widehat{c_p} \right  ) =
T \left ( \frac{p_0}{p_d} \right )^{R_d/\widehat{c_p}}
\left ( \frac{e}{e_s} \right )^{-r_v R_v/\widehat{c_p}}
\exp \left ( \frac{l_v\,r_v}{\widehat{c_p} T} \right )
\end{equation}
When $e=e_s$ \eqref{eq:truetheta} reduces to the $\theta_e$ we know and love, but
when the air is subsaturated, we can use \eqref{eq:truetheta} to calculate $\theta_e$
without having to bring the parcel up to LCL to ensure saturation.  Also note that,
unlike the pseudoadiabatic $\theta_e$ which discards liquid water, this version
is fully reversible and can be used anywhere we need to calculate an entropy.

Finally, you can take the differential of \eqref{eq:truetheta} and show that

\begin{equation}
  \label{eq:difftrue}
d\phi = \widehat{c_p} \frac{d \theta_e}{\theta_e}  = (c_{pd} + r_v c_{pv} +r_l c_l)
\frac{dT}{T} +c_l \log T d r_T + \frac{l_v d r_v}{T} - R_d \frac{dp}{p} 
\end{equation}
The form of $\theta_e$ given in \eqref{eq:truetheta} or \eqref{eq:difftrue}
is what you would need to the wet bulb problem exactly, but obviously
it's easier to go with $h_m$ for this, since all of the vapor information
is contained in $r_v$ in that equation, instead of being split
between $r_v= \epsilon e/(p-e)$ and $e$) as it is in \eqref{eq:truetheta}.

Question:  So why not use $h_m$ all the time?






\end{document}

%%% Local Variables:
%%% mode: latex
%%% TeX-master: t
%%% End:
